\documentclass[12pt,a4paper]{article}
\usepackage[utf8]{inputenc}
\usepackage[T1]{fontenc}
\usepackage[spanish]{babel}
\usepackage{amsmath,amssymb}
\usepackage{geometry}
\usepackage{siunitx}
\usepackage{enumitem}
\usepackage{graphicx}
\usepackage{booktabs}
\geometry{margin=2.5cm}

\title{Tarea --- Vectores Velocidad y Aceleración\\en Coordenadas Polares}
\author{Mecánica Aplicada}
\date{28 de mayo de 2026}

\begin{document}
\maketitle

\section{Introducción}

Las coordenadas polares $(r,\theta)$ son un sistema de coordenadas curvilíneas
ortogonales donde la posición de un punto $P$ se describe mediante su distancia
$r$ al origen y el ángulo $\theta$ que forma el radio vector con el eje polar.
Este sistema es particularmente útil cuando el movimiento tiene una componente
rotacional o cuando la trayectoria es curva.

Para describir el movimiento en coordenadas polares se definen dos vectores
unitarios:
\begin{itemize}
    \item $\mathbf{e}_r$: vector unitario radial, en dirección de $r$ creciente.
    \item $\mathbf{e}_\theta$: vector unitario transversal, perpendicular a
          $\mathbf{e}_r$ en dirección de $\theta$ creciente.
\end{itemize}

\section{Vectores unitarios y sus derivadas}

El vector de posición en coordenadas polares es:

\begin{equation}
    \mathbf{r} = r\,\mathbf{e}_r
    \label{eq:posicion}
\end{equation}

Los vectores unitarios $\mathbf{e}_r$ y $\mathbf{e}_\theta$ pueden expresarse
en coordenadas cartesianas como:

\begin{align}
    \mathbf{e}_r &= \cos\theta\,\mathbf{i} + \sin\theta\,\mathbf{j} \label{eq:er} \\
    \mathbf{e}_\theta &= -\sin\theta\,\mathbf{i} + \cos\theta\,\mathbf{j} \label{eq:etheta}
\end{align}

\subsection{Derivada de $\mathbf{e}_r$}

Derivando (\ref{eq:er}) respecto al tiempo:

\begin{align*}
    \dot{\mathbf{e}}_r &= \frac{d}{dt}(\cos\theta)\,\mathbf{i} + \frac{d}{dt}(\sin\theta)\,\mathbf{j} \\
    &= (-\sin\theta)\,\dot{\theta}\,\mathbf{i} + (\cos\theta)\,\dot{\theta}\,\mathbf{j} \\
    &= \dot{\theta}\,(-\sin\theta\,\mathbf{i} + \cos\theta\,\mathbf{j}) \\
    &= \dot{\theta}\,\mathbf{e}_\theta
\end{align*}

\begin{equation}
    \boxed{\dot{\mathbf{e}}_r = \dot{\theta}\,\mathbf{e}_\theta}
    \label{eq:der_er}
\end{equation}

\subsection{Derivada de $\mathbf{e}_\theta$}

Derivando (\ref{eq:etheta}) respecto al tiempo:

\begin{align*}
    \dot{\mathbf{e}}_\theta &= \frac{d}{dt}(-\sin\theta)\,\mathbf{i} + \frac{d}{dt}(\cos\theta)\,\mathbf{j} \\
    &= (-\cos\theta)\,\dot{\theta}\,\mathbf{i} + (-\sin\theta)\,\dot{\theta}\,\mathbf{j} \\
    &= -\dot{\theta}\,(\cos\theta\,\mathbf{i} + \sin\theta\,\mathbf{j}) \\
    &= -\dot{\theta}\,\mathbf{e}_r
\end{align*}

\begin{equation}
    \boxed{\dot{\mathbf{e}}_\theta = -\dot{\theta}\,\mathbf{e}_r}
    \label{eq:der_etheta}
\end{equation}

\section{Vector Velocidad}

La velocidad se define como la derivada temporal del vector de posición:

\[
\mathbf{v} = \dot{\mathbf{r}} = \frac{d}{dt}(r\,\mathbf{e}_r)
\]

Aplicando la regla del producto:

\begin{align*}
    \mathbf{v} &= \dot{r}\,\mathbf{e}_r + r\,\dot{\mathbf{e}}_r \\
    &= \dot{r}\,\mathbf{e}_r + r(\dot{\theta}\,\mathbf{e}_\theta) \\
    &= \dot{r}\,\mathbf{e}_r + r\dot{\theta}\,\mathbf{e}_\theta
\end{align*}

\begin{equation}
    \boxed{\mathbf{v} = \dot{r}\,\mathbf{e}_r + r\dot{\theta}\,\mathbf{e}_\theta}
    \label{eq:velocidad}
\end{equation}

\subsection{Componentes de la velocidad}

\begin{itemize}
    \item \textbf{Componente radial:} $v_r = \dot{r}$
    
    Representa la rapidez con la que cambia la distancia al origen.
    
    \item \textbf{Componente transversal:} $v_\theta = r\dot{\theta}$
    
    Representa la velocidad debida al cambio de dirección (movimiento angular).
    También se conoce como velocidad acimutal.
\end{itemize}

La magnitud de la velocidad es:

\[
|\mathbf{v}| = \sqrt{\dot{r}^2 + (r\dot{\theta})^2}
\]

\section{Vector Aceleración}

La aceleración se define como la derivada temporal del vector velocidad:

\[
\mathbf{a} = \dot{\mathbf{v}} = \frac{d}{dt}(\dot{r}\,\mathbf{e}_r + r\dot{\theta}\,\mathbf{e}_\theta)
\]

Aplicando la regla del producto a cada término:

\begin{align*}
    \mathbf{a} &= \frac{d}{dt}(\dot{r}\,\mathbf{e}_r) + \frac{d}{dt}(r\dot{\theta}\,\mathbf{e}_\theta) \\[4pt]
    &= \ddot{r}\,\mathbf{e}_r + \dot{r}\,\dot{\mathbf{e}}_r
       + \dot{r}\dot{\theta}\,\mathbf{e}_\theta + r\ddot{\theta}\,\mathbf{e}_\theta
       + r\dot{\theta}\,\dot{\mathbf{e}}_\theta
\end{align*}

Sustituyendo (\ref{eq:der_er}) y (\ref{eq:der_etheta}):

\begin{align*}
    \mathbf{a} &= \ddot{r}\,\mathbf{e}_r + \dot{r}(\dot{\theta}\,\mathbf{e}_\theta)
               + \dot{r}\dot{\theta}\,\mathbf{e}_\theta + r\ddot{\theta}\,\mathbf{e}_\theta
               + r\dot{\theta}(-\dot{\theta}\,\mathbf{e}_r) \\[4pt]
    &= \ddot{r}\,\mathbf{e}_r + \dot{r}\dot{\theta}\,\mathbf{e}_\theta
       + \dot{r}\dot{\theta}\,\mathbf{e}_\theta + r\ddot{\theta}\,\mathbf{e}_\theta
       - r\dot{\theta}^2\,\mathbf{e}_r \\[4pt]
    &= (\ddot{r} - r\dot{\theta}^2)\,\mathbf{e}_r
       + (2\dot{r}\dot{\theta} + r\ddot{\theta})\,\mathbf{e}_\theta
\end{align*}

\begin{equation}
    \boxed{\mathbf{a} = (\ddot{r} - r\dot{\theta}^2)\,\mathbf{e}_r
           + (2\dot{r}\dot{\theta} + r\ddot{\theta})\,\mathbf{e}_\theta}
    \label{eq:aceleracion}
\end{equation}

\subsection{Componentes de la aceleración}

\begin{itemize}
    \item \textbf{Componente radial:} $a_r = \ddot{r} - r\dot{\theta}^2$
    
    El término $\ddot{r}$ es la aceleración radial debida al cambio en la
    distancia al origen. El término $-r\dot{\theta}^2$ es la aceleración
    centrípeta, dirigida hacia el origen.
    
    \item \textbf{Componente transversal:} $a_\theta = 2\dot{r}\dot{\theta} + r\ddot{\theta}$
    
    El término $2\dot{r}\dot{\theta}$ es la aceleración de Coriolis (en el
    contexto de coordenadas polares). El término $r\ddot{\theta}$ es la
    aceleración angular.
\end{itemize}

La magnitud de la aceleración es:

\[
|\mathbf{a}| = \sqrt{(\ddot{r} - r\dot{\theta}^2)^2 + (2\dot{r}\dot{\theta} + r\ddot{\theta})^2}
\]

\section{Forma alternativa de la componente transversal}

La componente transversal puede expresarse de forma más compacta:

\begin{align*}
    a_\theta &= 2\dot{r}\dot{\theta} + r\ddot{\theta} \\
    &= \frac{1}{r}\frac{d}{dt}(r^2\dot{\theta})
\end{align*}

Demostración:

\[
\frac{d}{dt}(r^2\dot{\theta}) = 2r\dot{r}\dot{\theta} + r^2\ddot{\theta}
   = r(2\dot{r}\dot{\theta} + r\ddot{\theta}) = r\,a_\theta
\]

\[
\therefore a_\theta = \frac{1}{r}\frac{d}{dt}(r^2\dot{\theta})
\]

Esta forma es útil cuando se conserva el momento angular ($r^2\dot{\theta} = \text{cte}$).

\section{Casos particulares}

\subsection{Movimiento circular ($r = \text{cte}$, $\dot{r} = 0$, $\ddot{r} = 0$)}

\begin{align*}
    \mathbf{v} &= r\dot{\theta}\,\mathbf{e}_\theta \\
    \mathbf{a} &= -r\dot{\theta}^2\,\mathbf{e}_r + r\ddot{\theta}\,\mathbf{e}_\theta
\end{align*}

Para movimiento circular uniforme ($\ddot{\theta} = 0$):
\[
\mathbf{a} = -r\dot{\theta}^2\,\mathbf{e}_r = -\frac{v^2}{r}\,\mathbf{e}_r
\]

\subsection{Movimiento rectilíneo radial ($\dot{\theta} = 0$, $\ddot{\theta} = 0$)}

\begin{align*}
    \mathbf{v} &= \dot{r}\,\mathbf{e}_r \\
    \mathbf{a} &= \ddot{r}\,\mathbf{e}_r
\end{align*}

\section{Resumen de expresiones}

\begin{center}
\begin{tabular}{lll}
\toprule
Magnitud & Expresión & Componentes \\
\midrule
Posición & $\mathbf{r} = r\,\mathbf{e}_r$ & --- \\
Velocidad & $\mathbf{v} = \dot{r}\,\mathbf{e}_r + r\dot{\theta}\,\mathbf{e}_\theta$ &
$v_r = \dot{r},\; v_\theta = r\dot{\theta}$ \\
Aceleración & $\mathbf{a} = (\ddot{r} - r\dot{\theta}^2)\,\mathbf{e}_r
            + (2\dot{r}\dot{\theta} + r\ddot{\theta})\,\mathbf{e}_\theta$ &
$a_r = \ddot{r} - r\dot{\theta}^2,\; a_\theta = 2\dot{r}\dot{\theta} + r\ddot{\theta}$ \\
\bottomrule
\end{tabular}
\end{center}

\section*{Derivadas de vectores unitarios (resumen)}

\[
\boxed{\dot{\mathbf{e}}_r = \dot{\theta}\,\mathbf{e}_\theta}
\qquad
\boxed{\dot{\mathbf{e}}_\theta = -\dot{\theta}\,\mathbf{e}_r}
\]

\end{document}
